% Erratum E1 for inclusion in main.tex.
% Suggested use: % Erratum E1 for inclusion in main.tex.
% Suggested use: % Erratum E1 for inclusion in main.tex.
% Suggested use: % Erratum E1 for inclusion in main.tex.
% Suggested use: \input{erratum_E1_Ncomm_endpoint}

\subsection*{Erratum E1: the endpoint formulation of equation (6)}
\label{s:erratum:E1}

In the sentence preceding equation (6), the range
$0 \leq \theta \leq 1$ should be replaced by
$0 < \theta < 1$ if the displayed value of $\alpha$ is used.  Indeed,
\[
	\alpha
	=
	t^{1/k}
	\left(\frac{1-\theta}{\theta}\right)^{1-1/(2k)}
\]
is not defined at the endpoints $\theta=0$ and $\theta=1$.
The endpoint cases should instead be treated separately, or obtained by
an explicit limiting argument.

For completeness, let us recall the verification of equation (6) in
the corrected range $0<\theta<1$.  Put
\[
	p := 2k \geq 1 .
\]
With the displayed choice of $\alpha$, one has
\[
	\alpha^p
	=
	t^2
	\left(\frac{1-\theta}{\theta}\right)^{p-1},
	\qquad\text{or equivalently}\qquad
	t^2
	=
	\alpha^p
	\left(\frac{\theta}{1-\theta}\right)^{p-1} .
\]
Using $N_\theta^{-2}=\theta^\theta(1-\theta)^{1-\theta}$, the left hand
side of equation (6) becomes
\[
	N_\theta^{2(1-p)}
	\frac{t^{2(1-\theta)}}{t^2+\lambda^p}
	=
	\frac{
		\theta^{p-1}\alpha^{p(1-\theta)}
	}{
		\left(\frac{\theta}{1-\theta}\right)^{p-1}\alpha^p
		+
		\lambda^p
	} .
\]
Thus equation (6) is equivalent to
\[
	(\alpha+\lambda)^p
	\leq
	(1-\theta)^{1-p}\alpha^p
	+
	\theta^{1-p}\lambda^p .
\]
This last estimate follows from convexity of $s\mapsto s^p$:
\[
	(\alpha+\lambda)^p
	=
	\left(
		(1-\theta)\frac{\alpha}{1-\theta}
		+
		\theta\frac{\lambda}{\theta}
	\right)^p
	\leq
	(1-\theta)^{1-p}\alpha^p
	+
	\theta^{1-p}\lambda^p .
\]
This proves the corrected form of equation (6).

At the endpoints the displayed value of $\alpha$ should not be used.
The relevant limiting scalar estimates are elementary.  For $\theta=0$,
$N_0=1$ and
\[
	\frac{t^2}{t^2+\lambda^{2k}}
	\leq
	1
	=
	\lim_{\alpha\to\infty}
	\left(\frac{\alpha}{\alpha+\lambda}\right)^{2k} .
\]
For $\theta=1$, $N_1=1$ and, for $\lambda>0$,
\[
	\frac{1}{t^2+\lambda^{2k}}
	\leq
	\lambda^{-2k}
	=
	\lim_{\alpha\downarrow0}
	\left(\frac{1}{\alpha+\lambda}\right)^{2k} .
\]
The case $\lambda=0$ is interpreted in the extended sense in the last
inequality; in the spectral applications of the paper, mass at
$\lambda=0$ is absent because $T$ is assumed injective.  Hence any proof
which uses equation (6) at $\theta=0$ or $\theta=1$ should replace
that invocation by the corresponding endpoint argument.


\subsection*{Erratum E1: the endpoint formulation of equation (6)}
\label{s:erratum:E1}

In the sentence preceding equation (6), the range
$0 \leq \theta \leq 1$ should be replaced by
$0 < \theta < 1$ if the displayed value of $\alpha$ is used.  Indeed,
\[
	\alpha
	=
	t^{1/k}
	\left(\frac{1-\theta}{\theta}\right)^{1-1/(2k)}
\]
is not defined at the endpoints $\theta=0$ and $\theta=1$.
The endpoint cases should instead be treated separately, or obtained by
an explicit limiting argument.

For completeness, let us recall the verification of equation (6) in
the corrected range $0<\theta<1$.  Put
\[
	p := 2k \geq 1 .
\]
With the displayed choice of $\alpha$, one has
\[
	\alpha^p
	=
	t^2
	\left(\frac{1-\theta}{\theta}\right)^{p-1},
	\qquad\text{or equivalently}\qquad
	t^2
	=
	\alpha^p
	\left(\frac{\theta}{1-\theta}\right)^{p-1} .
\]
Using $N_\theta^{-2}=\theta^\theta(1-\theta)^{1-\theta}$, the left hand
side of equation (6) becomes
\[
	N_\theta^{2(1-p)}
	\frac{t^{2(1-\theta)}}{t^2+\lambda^p}
	=
	\frac{
		\theta^{p-1}\alpha^{p(1-\theta)}
	}{
		\left(\frac{\theta}{1-\theta}\right)^{p-1}\alpha^p
		+
		\lambda^p
	} .
\]
Thus equation (6) is equivalent to
\[
	(\alpha+\lambda)^p
	\leq
	(1-\theta)^{1-p}\alpha^p
	+
	\theta^{1-p}\lambda^p .
\]
This last estimate follows from convexity of $s\mapsto s^p$:
\[
	(\alpha+\lambda)^p
	=
	\left(
		(1-\theta)\frac{\alpha}{1-\theta}
		+
		\theta\frac{\lambda}{\theta}
	\right)^p
	\leq
	(1-\theta)^{1-p}\alpha^p
	+
	\theta^{1-p}\lambda^p .
\]
This proves the corrected form of equation (6).

At the endpoints the displayed value of $\alpha$ should not be used.
The relevant limiting scalar estimates are elementary.  For $\theta=0$,
$N_0=1$ and
\[
	\frac{t^2}{t^2+\lambda^{2k}}
	\leq
	1
	=
	\lim_{\alpha\to\infty}
	\left(\frac{\alpha}{\alpha+\lambda}\right)^{2k} .
\]
For $\theta=1$, $N_1=1$ and, for $\lambda>0$,
\[
	\frac{1}{t^2+\lambda^{2k}}
	\leq
	\lambda^{-2k}
	=
	\lim_{\alpha\downarrow0}
	\left(\frac{1}{\alpha+\lambda}\right)^{2k} .
\]
The case $\lambda=0$ is interpreted in the extended sense in the last
inequality; in the spectral applications of the paper, mass at
$\lambda=0$ is absent because $T$ is assumed injective.  Hence any proof
which uses equation (6) at $\theta=0$ or $\theta=1$ should replace
that invocation by the corresponding endpoint argument.


\subsection*{Erratum E1: the endpoint formulation of equation (6)}
\label{s:erratum:E1}

In the sentence preceding equation (6), the range
$0 \leq \theta \leq 1$ should be replaced by
$0 < \theta < 1$ if the displayed value of $\alpha$ is used.  Indeed,
\[
	\alpha
	=
	t^{1/k}
	\left(\frac{1-\theta}{\theta}\right)^{1-1/(2k)}
\]
is not defined at the endpoints $\theta=0$ and $\theta=1$.
The endpoint cases should instead be treated separately, or obtained by
an explicit limiting argument.

For completeness, let us recall the verification of equation (6) in
the corrected range $0<\theta<1$.  Put
\[
	p := 2k \geq 1 .
\]
With the displayed choice of $\alpha$, one has
\[
	\alpha^p
	=
	t^2
	\left(\frac{1-\theta}{\theta}\right)^{p-1},
	\qquad\text{or equivalently}\qquad
	t^2
	=
	\alpha^p
	\left(\frac{\theta}{1-\theta}\right)^{p-1} .
\]
Using $N_\theta^{-2}=\theta^\theta(1-\theta)^{1-\theta}$, the left hand
side of equation (6) becomes
\[
	N_\theta^{2(1-p)}
	\frac{t^{2(1-\theta)}}{t^2+\lambda^p}
	=
	\frac{
		\theta^{p-1}\alpha^{p(1-\theta)}
	}{
		\left(\frac{\theta}{1-\theta}\right)^{p-1}\alpha^p
		+
		\lambda^p
	} .
\]
Thus equation (6) is equivalent to
\[
	(\alpha+\lambda)^p
	\leq
	(1-\theta)^{1-p}\alpha^p
	+
	\theta^{1-p}\lambda^p .
\]
This last estimate follows from convexity of $s\mapsto s^p$:
\[
	(\alpha+\lambda)^p
	=
	\left(
		(1-\theta)\frac{\alpha}{1-\theta}
		+
		\theta\frac{\lambda}{\theta}
	\right)^p
	\leq
	(1-\theta)^{1-p}\alpha^p
	+
	\theta^{1-p}\lambda^p .
\]
This proves the corrected form of equation (6).

At the endpoints the displayed value of $\alpha$ should not be used.
The relevant limiting scalar estimates are elementary.  For $\theta=0$,
$N_0=1$ and
\[
	\frac{t^2}{t^2+\lambda^{2k}}
	\leq
	1
	=
	\lim_{\alpha\to\infty}
	\left(\frac{\alpha}{\alpha+\lambda}\right)^{2k} .
\]
For $\theta=1$, $N_1=1$ and, for $\lambda>0$,
\[
	\frac{1}{t^2+\lambda^{2k}}
	\leq
	\lambda^{-2k}
	=
	\lim_{\alpha\downarrow0}
	\left(\frac{1}{\alpha+\lambda}\right)^{2k} .
\]
The case $\lambda=0$ is interpreted in the extended sense in the last
inequality; in the spectral applications of the paper, mass at
$\lambda=0$ is absent because $T$ is assumed injective.  Hence any proof
which uses equation (6) at $\theta=0$ or $\theta=1$ should replace
that invocation by the corresponding endpoint argument.


\subsection*{Erratum E1: the endpoint formulation of equation (6)}
\label{s:erratum:E1}

In the sentence preceding equation (6), the range
$0 \leq \theta \leq 1$ should be replaced by
$0 < \theta < 1$ if the displayed value of $\alpha$ is used.  Indeed,
\[
	\alpha
	=
	t^{1/k}
	\left(\frac{1-\theta}{\theta}\right)^{1-1/(2k)}
\]
is not defined at the endpoints $\theta=0$ and $\theta=1$.
The endpoint cases should instead be treated separately, or obtained by
an explicit limiting argument.

For completeness, let us recall the verification of equation (6) in
the corrected range $0<\theta<1$.  Put
\[
	p := 2k \geq 1 .
\]
With the displayed choice of $\alpha$, one has
\[
	\alpha^p
	=
	t^2
	\left(\frac{1-\theta}{\theta}\right)^{p-1},
	\qquad\text{or equivalently}\qquad
	t^2
	=
	\alpha^p
	\left(\frac{\theta}{1-\theta}\right)^{p-1} .
\]
Using $N_\theta^{-2}=\theta^\theta(1-\theta)^{1-\theta}$, the left hand
side of equation (6) becomes
\[
	N_\theta^{2(1-p)}
	\frac{t^{2(1-\theta)}}{t^2+\lambda^p}
	=
	\frac{
		\theta^{p-1}\alpha^{p(1-\theta)}
	}{
		\left(\frac{\theta}{1-\theta}\right)^{p-1}\alpha^p
		+
		\lambda^p
	} .
\]
Thus equation (6) is equivalent to
\[
	(\alpha+\lambda)^p
	\leq
	(1-\theta)^{1-p}\alpha^p
	+
	\theta^{1-p}\lambda^p .
\]
This last estimate follows from convexity of $s\mapsto s^p$:
\[
	(\alpha+\lambda)^p
	=
	\left(
		(1-\theta)\frac{\alpha}{1-\theta}
		+
		\theta\frac{\lambda}{\theta}
	\right)^p
	\leq
	(1-\theta)^{1-p}\alpha^p
	+
	\theta^{1-p}\lambda^p .
\]
This proves the corrected form of equation (6).

At the endpoints the displayed value of $\alpha$ should not be used.
The relevant limiting scalar estimates are elementary.  For $\theta=0$,
$N_0=1$ and
\[
	\frac{t^2}{t^2+\lambda^{2k}}
	\leq
	1
	=
	\lim_{\alpha\to\infty}
	\left(\frac{\alpha}{\alpha+\lambda}\right)^{2k} .
\]
For $\theta=1$, $N_1=1$ and, for $\lambda>0$,
\[
	\frac{1}{t^2+\lambda^{2k}}
	\leq
	\lambda^{-2k}
	=
	\lim_{\alpha\downarrow0}
	\left(\frac{1}{\alpha+\lambda}\right)^{2k} .
\]
The case $\lambda=0$ is interpreted in the extended sense in the last
inequality; in the spectral applications of the paper, mass at
$\lambda=0$ is absent because $T$ is assumed injective.  Hence any proof
which uses equation (6) at $\theta=0$ or $\theta=1$ should replace
that invocation by the corresponding endpoint argument.
